% =========================================================
% Monotone Functions and Continuity — Canonical Notes
% Chapter: continuity
% Volume III · Analysis
% =========================================================

\section{Monotone Functions and Continuity}
\label{sec:monotone-functions}

% ---------------------------------------------------------
% Toolkit
% ---------------------------------------------------------

\begin{tcolorbox}[title={Toolkit — Monotone Functions}]
\begin{itemize}
  \item Increasing $f$ has left and right limits at every interior point.
  \item $f$ continuous at $c$ $\iff$ $f(c^-) = f(c) = f(c^+)$ $\iff$ jump $j_f(c) = 0$.
  \item Every discontinuity of a monotone function is of first kind (jump type).
  \item Jump intervals at distinct discontinuities are disjoint.
  \item Discontinuities of a monotone function are countable.
  \item Strictly monotone continuous bijection on $[a,b]$ has a continuous inverse.
  \item Continuous injective on $[a,b]$ $\iff$ strictly monotone.
\end{itemize}
\end{tcolorbox}

% ---------------------------------------------------------
% Theorem — One-Sided Limits of Monotone Functions
% ---------------------------------------------------------

\begin{theorem}[One-Sided Limits of Monotone Functions]
\label{thm:monotone-one-sided-limits}
Let $I \subseteq \mathbb{R}$ be an interval and let
$f : I \to \mathbb{R}$ be increasing on $I$. Suppose $c \in I$ is
not an endpoint of $I$. Then:
\begin{enumerate}[label=\textup{(\roman*)}]
  \item $\displaystyle\lim_{x \to c^-} f(x) = \sup\{f(x) : x \in I,\; x < c\}$,
  \item $\displaystyle\lim_{x \to c^+} f(x) = \inf\{f(x) : x \in I,\; x > c\}$.
\end{enumerate}
In particular, both one-sided limits exist and are finite at every
interior point of $I$.

\smallskip
\noindent
\hyperref[prf:monotone-one-sided-limits]{\textit{Go to proof.}}
\end{theorem}

\begin{remark*}[Standard quantified statement]
\[
  f \text{ increasing on } I,\;
  c \text{ interior to } I
  \;\Rightarrow\;
  \lim_{x \to c^-} f = \sup\{f(x): x < c\}
  \;\wedge\;
  \lim_{x \to c^+} f = \inf\{f(x): x > c\}.
\]
\end{remark*}

\begin{remark*}[Interpretation]
Monotonicity guarantees that the left and right limits exist at every
interior point, even at discontinuities. The left limit is the
supremum of the function values to the left; the right limit is the
infimum of values to the right. The proof uses the $\varepsilon$
characterization of supremum to produce the required $\delta$.
\end{remark*}

% ---------------------------------------------------------
% Corollary — Continuity Criterion for Monotone Functions
% ---------------------------------------------------------

\begin{corollary}[Continuity Criterion for Monotone Functions]
\label{cor:monotone-continuity-criterion}
Let $I \subseteq \mathbb{R}$ be an interval and let
$f : I \to \mathbb{R}$ be increasing on $I$. Suppose $c \in I$ is
not an endpoint. The following are equivalent:
\begin{enumerate}[label=\textup{(\roman*)}]
  \item $f$ is continuous at $c$.
  \item $\displaystyle\lim_{x \to c^-} f(c) = f(c) = \lim_{x \to c^+} f(c)$.
  \item $\sup\{f(x): x < c\} = f(c) = \inf\{f(x): x > c\}$.
\end{enumerate}

\smallskip
\noindent
\hyperref[prf:monotone-continuity-criterion]{\textit{Go to proof.}}
\end{corollary>

\begin{remark*}[Standard quantified statement]
\[
  f \text{ continuous at } c
  \;\iff\;
  f(c^-) = f(c) = f(c^+)
  \;\iff\;
  \sup\{f(x): x < c\} = f(c) = \inf\{f(x): x > c\}.
\]
\end{remark*}

\begin{remark*}[Interpretation]
Continuity at $c$ is equivalent to the absence of a gap between the
left supremum and the right infimum at $c$. If either one-sided limit
differs from $f(c)$, there is a jump discontinuity.
\end{remark*}

% ---------------------------------------------------------
% Definition — Jump of a Function
% ---------------------------------------------------------

\begin{tcolorbox}[title={Definition — Jump of a Function}]
\begin{definition}[Jump of a Function]
\label{def:jump-of-function}
Let $I \subseteq \mathbb{R}$ be an interval and let
$f : I \to \mathbb{R}$ be increasing on $I$. The \textbf{jump} of $f$
at $c \in I$ is
\[
  j_f(c) \;:=\; \lim_{x \to c^+} f(x) - \lim_{x \to c^-} f(x).
\]
For an increasing function, $f$ is continuous at $c$ if and only if
$j_f(c) = 0$.
\end{definition}
\end{tcolorbox>

\begin{remark*}[Standard quantified statement]
\[
  j_f(c) = f(c^+) - f(c^-).
  \qquad
  f \text{ continuous at } c \;\iff\; j_f(c) = 0.
\]
\end{remark*}

\begin{remark*}[Interpretation]
The jump $j_f(c) \geq 0$ measures the gap in the graph at $c$.
Positive jump indicates a jump discontinuity; zero jump indicates
continuity. The graph of $f$ has a vertical gap of height $j_f(c)$
at each discontinuity.
\end{remark*}

% ---------------------------------------------------------
% Proposition — Discontinuities of a Monotone Function Are of First Kind
% ---------------------------------------------------------

\begin{proposition}[Discontinuities of a Monotone Function Are of First Kind]
\label{prop:monotone-discontinuities-first-kind}
Let $f : E \to \mathbb{R}$ be monotone on $E \subseteq \mathbb{R}$.
Then every discontinuity of $f$ is of first kind: at every
discontinuity $a \in E$, both one-sided limits $f(a^-)$ and $f(a^+)$
exist.

\smallskip
\noindent
\hyperref[prf:monotone-discontinuities-first-kind]{\textit{Go to proof.}}
\end{proposition>

\begin{remark*}[Standard quantified statement]
\[
  f \text{ monotone on } E,\;
  a \in E \text{ a discontinuity}
  \;\Rightarrow\;
  \lim_{x \to a^-} f(x) \in \mathbb{R}
  \;\wedge\;
  \lim_{x \to a^+} f(x) \in \mathbb{R}.
\]
\end{remark*}

\begin{remark*}[Interpretation]
A monotone function may jump but cannot oscillate. Oscillation would
require the function to move up and down near the discontinuity,
contradicting monotonicity. Hence all discontinuities are jump-type
(first kind). This is a strong structural constraint: monotonicity
rules out infinite and oscillatory discontinuities entirely.
\end{remark*}

% ---------------------------------------------------------
% Corollary — Jump Intervals of a Monotone Function Are Disjoint
% ---------------------------------------------------------

\begin{corollary}[Jump Intervals of a Monotone Function Are Disjoint]
\label{cor:jump-intervals-for-monotone-discontinuities}
Let $f : E \to \mathbb{R}$ be nondecreasing on $E$ and let $a \in E$
be a discontinuity. The \textbf{jump interval} at $a$ is the open
interval $(f(a^-), f(a^+))$. The function $f$ assumes no values in
this interval. Jump intervals at distinct discontinuities are pairwise
disjoint.

\smallskip
\noindent
\hyperref[prf:jump-intervals-for-monotone-discontinuities]{\textit{Go to proof.}}
\end{corollary>

\begin{remark*}[Standard quantified statement]
\[
  a \neq b \text{ both discontinuities of nondecreasing } f
  \;\Rightarrow\;
  (f(a^-), f(a^+)) \cap (f(b^-), f(b^+)) = \emptyset.
\]
\end{remark*}

\begin{remark*}[Interpretation]
If $a < b$ are both discontinuities of a nondecreasing $f$, then for
any $x \in (a,b)$, $f(a^+) \leq f(x) \leq f(b^-)$. The jump
interval at $a$ lies entirely below $f(b^-)$ and the jump interval
at $b$ lies entirely above $f(a^+)$, so they cannot overlap. This
disjointness is the geometric heart of the countability theorem.
\end{remark*}

% ---------------------------------------------------------
% Proposition — Criterion for Continuity of a Monotone Function
% ---------------------------------------------------------

\begin{proposition}[Criterion for Continuity of a Monotone Function]
\label{prop:criterion-for-continuity-of-monotone-function}
A monotone function $f : [a,b] \to \mathbb{R}$ is continuous on
$[a,b]$ if and only if its image $f([a,b])$ is the closed interval
with endpoints $f(a)$ and $f(b)$.

\smallskip
\noindent
\hyperref[prf:criterion-continuity-monotone]{\textit{Go to proof.}}
\end{proposition>

\begin{remark*}[Standard quantified statement]
\[
  f \text{ monotone on } [a,b] :\;
  f \text{ continuous on } [a,b]
  \;\iff\;
  f([a,b]) = \bigl[\min(f(a),f(b)),\; \max(f(a),f(b))\bigr].
\]
\end{remark*}

\begin{remark*}[Interpretation]
For a monotone function, continuity is equivalent to having no gaps
in the range. Each jump discontinuity creates a gap (its jump
interval); continuity is exactly the condition that no such gaps
exist. This is the converse of the Image Theorem for monotone
functions: having image $[f(a),f(b)]$ forces continuity.
\end{remark*}

% ---------------------------------------------------------
% Theorem — Discontinuities of a Monotone Function Are Countable
% ---------------------------------------------------------

\begin{theorem}[Discontinuities of a Monotone Function Are Countable]
\label{thm:monotone-discontinuities-countable}
Let $I \subseteq \mathbb{R}$ be an interval and let
$f : I \to \mathbb{R}$ be monotone on $I$. Then the set
$D := \{c \in I : f \text{ discontinuous at } c\}$ is countable.

\smallskip
\noindent
\hyperref[prf:monotone-discontinuities-countable]{\textit{Go to proof.}}
\end{theorem}

\begin{remark*}[Standard quantified statement]
\[
  f \text{ monotone on } I
  \;\Rightarrow\;
  D \text{ is countable}.
\]
\end{remark*}

\begin{remark*}[Interpretation]
At each discontinuity $c$, the jump $j_f(c) > 0$ and the jump
interval $(f(c^-), f(c^+))$ is nonempty. These intervals are
pairwise disjoint. Each nonempty open interval contains a rational,
so the map $c \mapsto$ (a rational in $(f(c^-), f(c^+))$) is an
injection from $D$ into $\mathbb{Q}$. Hence $D$ is countable.

This is a sharp result: monotone functions can have countably
infinitely many discontinuities, but never uncountably many. The
proof uses density of $\mathbb{Q}$ as the key tool.
\end{remark*>

% ---------------------------------------------------------
% Corollary — Monotone Function Is Continuous on a Dense Set
% ---------------------------------------------------------

\begin{corollary}[Monotone Function Is Continuous on a Dense Set]
\label{cor:monotone-continuous-on-dense-set}
Let $f$ be monotone on $[a,b]$. Then $f$ is continuous at the points
of a set dense in $[a,b]$.

\smallskip
\noindent
\hyperref[prf:monotone-continuous-on-dense-set]{\textit{Go to proof.}}
\end{corollary>

\begin{remark*}[Standard quantified statement]
\[
  f \text{ monotone on } [a,b]
  \;\Rightarrow\;
  C_f := \{x \in [a,b] : f \text{ continuous at } x\}
  \text{ is dense in } [a,b].
\]
\end{remark*>

\begin{remark*}[Interpretation]
The discontinuity set $D_f$ is countable, hence nowhere dense. Its
complement $C_f = [a,b] \setminus D_f$ is therefore dense. Regardless
of how many jump discontinuities a monotone function has, they can
never fill an interval; the function is well-behaved at most points
in any neighborhood.
\end{remark*>

% ---------------------------------------------------------
% Proposition — Inverse of a Strictly Monotone Function
% ---------------------------------------------------------

\begin{proposition}[Inverse of a Strictly Monotone Function]
\label{prop:inverse-of-strictly-monotone-function}
Let $f : X \to \mathbb{R}$ be strictly monotone on
$X \subseteq \mathbb{R}$. Then $f$ has an inverse
$f^{-1} : f(X) \to X$, and $f^{-1}$ has the same type of
monotonicity on $f(X)$ as $f$ has on $X$.
\end{proposition>

\begin{remark*}[Standard quantified statement]
\[
  f \text{ strictly increasing on } X
  \;\Rightarrow\;
  f^{-1} \text{ strictly increasing on } f(X).
\]
Similarly for strictly decreasing.
\end{remark*>

\begin{remark*}[Interpretation]
Strict monotonicity implies injectivity, so the inverse exists.
The order direction is preserved: if $y_1 < y_2$ are in $f(X)$,
say $y_i = f(x_i)$, then $f^{-1}(y_1) = x_1 < x_2 = f^{-1}(y_2)$
because $f$ is strictly increasing.
\end{remark*>

% ---------------------------------------------------------
% Proposition — Continuous Injective Is Strictly Monotone
% ---------------------------------------------------------

\begin{proposition}[Continuous Injective Map on a Closed Interval Is Strictly Monotone]
\label{prop:continuous-injective-iff-strictly-monotone}
A continuous mapping $f : [a,b] \to \mathbb{R}$ is injective if and
only if $f$ is strictly monotone on $[a,b]$.

\smallskip
\noindent
\hyperref[prf:continuous-injective-is-strictly-monotone]{\textit{Go to proof.}}
\end{proposition>

\begin{remark*}[Standard quantified statement]
\[
  f \text{ continuous on } [a,b] :\;
  f \text{ injective}
  \;\iff\;
  f \text{ strictly increasing or strictly decreasing on } [a,b].
\]
\end{remark*>

\begin{remark*}[Interpretation]
The forward direction is the nontrivial one: non-monotonicity forces
a local extremum, and IVT then produces two distinct preimages of
some value, contradicting injectivity. Strict monotonicity trivially
implies injectivity.
\end{remark*>

% ---------------------------------------------------------
% Theorem — Inverse Function Theorem for Monotone Functions
% ---------------------------------------------------------

\begin{theorem}[Inverse Function Theorem for Continuous Strictly Monotone Functions]
\label{thm:inverse-function-theorem-monotone}
Let $f : X \to \mathbb{R}$ be strictly monotone on $X \subseteq
\mathbb{R}$. Then $f^{-1} : f(X) \to X$ exists and has the same
monotonicity type as $f$.

If in addition $X = [a,b]$ and $f$ is continuous on $[a,b]$, then
$f([a,b])$ is the closed interval with endpoints $f(a)$ and $f(b)$,
and $f^{-1}$ is continuous on $f([a,b])$.
\end{theorem>

\begin{remark*}[Standard quantified statement]
\[
  f \text{ strictly monotone and continuous on } [a,b]
  \;\Rightarrow\;
  f([a,b]) = [\min(f(a),f(b)),\, \max(f(a),f(b))]
  \;\wedge\;
  f^{-1} \text{ continuous on } f([a,b]).
\]
\end{remark*>

\begin{remark*}[Interpretation]
Three conclusions: (1) the inverse exists and inherits monotonicity;
(2) the image is a closed interval; (3) the inverse is continuous.
The image identification uses the Criterion for Continuity of a
Monotone Function. Continuity of the inverse follows because a jump
in $f^{-1}$ would create a gap in $f([a,b])$, contradicting
that $f([a,b])$ is an interval.
\end{remark*>

% ---------------------------------------------------------
% Theorem — Continuous Inverse Theorem
% ---------------------------------------------------------

\begin{theorem}[Continuous Inverse Theorem]
\label{thm:continuous-inverse-theorem}
Let $I \subseteq \mathbb{R}$ be an interval and let
$f : I \to \mathbb{R}$ be strictly monotone and continuous on $I$.
Then the inverse $g := f^{-1}$ is strictly monotone and continuous
on $J := f(I)$.

\smallskip
\noindent
\hyperref[prf:inverse-strictly-monotone]{\textit{Go to proof.}}
\end{theorem>

\begin{remark*}[Standard quantified statement]
\[
  f \text{ strictly monotone and continuous on interval } I
  \;\Rightarrow\;
  f^{-1} : f(I) \to I \text{ strictly monotone and continuous on } f(I).
\]
\end{remark*>

\begin{remark*}[Interpretation]
Continuity of $f$ and Preservation of Intervals ensure $f(I)$ is an
interval. If $f^{-1}$ had a jump discontinuity at some $d \in J$,
then the jump interval around $d$ would be a gap in $f(I)$,
contradicting that $f(I)$ is an interval. Hence $f^{-1}$ is
continuous.
\end{remark*>

% ---------------------------------------------------------
% Theorem — Injective Continuous Function Is Strictly Monotone
% ---------------------------------------------------------

\begin{theorem}[Injective Continuous Function Is Strictly Monotone]
\label{thm:injective-continuous-is-monotone}
Let $I \subseteq \mathbb{R}$ be an interval and let
$f : I \to \mathbb{R}$ be injective and continuous on $I$. Then $f$
is strictly monotone on $I$.

\smallskip
\noindent
\hyperref[prf:continuous-injective-is-strictly-monotone]{\textit{Go to proof.}}
\end{theorem>

\begin{remark*}[Standard quantified statement]
\[
  f \text{ injective and continuous on interval } I
  \;\Rightarrow\;
  f \text{ strictly increasing or strictly decreasing on } I.
\]
\end{remark*>

\begin{remark*}[Interpretation]
If $f$ is not monotone, there exist $a < b < c$ in $I$ with a local
extremum at $b$. IVT applied to the appropriate subintervals
produces two distinct preimages of some intermediate value,
contradicting injectivity. The proof is a formal exploitation of
the order structure of $\mathbb{R}$.
\end{remark*>
